% preamble/cover.tex
% ============================================
%  封面设计 —— 墨绿 + 暗金配色
%  基于大笔记模板，适配常微分方程笔记
%
%  接口（在 main.tex 中使用）：
%    \booktitle{标题}        —— 封面标题
%    \bookauthor{作者}       —— 作者
%    \bookdate{日期}         —— 日期
%    \seriestitle{系列名}     —— 系列标识
%    \coverimage{封面图片路径} —— 底部封面大图
%    \makecover              —— 生成封面
% ============================================

\makeatletter

% ---- 可配置的封面元信息 ----
\newcommand{\@seriestitle}{Notes in Mathematics}
\newcommand{\seriestitle}[1]{\renewcommand{\@seriestitle}{#1}}

\newcommand{\@booktitle}{笔记标题}
\newcommand{\booktitle}[1]{\renewcommand{\@booktitle}{#1}}

\newcommand{\@bookauthor}{作者姓名}
\newcommand{\bookauthor}[1]{\renewcommand{\@bookauthor}{#1}}

\newcommand{\@booksubtitle}{ORDINARY DIFFERENTIAL EQUATIONS \& DYNAMICAL SYSTEMS}
\newcommand{\booksubtitle}[1]{\renewcommand{\@booksubtitle}{#1}}

\newcommand{\@bookdate}{\today}
\newcommand{\bookdate}[1]{\renewcommand{\@bookdate}{#1}}

\providecommand{\@coverimage}{}
\newcommand{\coverimage}[1]{\renewcommand{\@coverimage}{#1}}

% ---- 生成封面 ----
\newcommand{\makecover}{%
\begin{titlepage}
\thispagestyle{empty}

% ============================================================
%  第一层：左侧深蓝色竖条 + 渐变纹理 + 装饰圆点
% ============================================================
\begin{tikzpicture}[remember picture,overlay]
  \fill[BlueMain!90]
    (current page.south west) rectangle
    ([xshift=2.8cm]current page.north west);
  \shade[top color=BlueDark!90, bottom color=BlueMain!90]
    ([xshift=0cm,yshift=-5cm]current page.north west) rectangle
    ([xshift=2.8cm]current page.north west);
  \shade[bottom color=BlueDark!90, top color=BlueMain!90]
    ([xshift=0cm,yshift=5cm]current page.south west) rectangle
    ([xshift=2.8cm]current page.south west);
  \draw[BlueLight!40,line width=1pt]
    ([xshift=2.8cm]current page.south west) --
    ([xshift=2.8cm]current page.north west);
  \draw[BlueLight!20,line width=0.4pt]
    ([xshift=0.4cm]current page.south west) --
    ([xshift=0.4cm]current page.north west);
  \foreach \y/\r/\o in {
    0.08/1.1/0.06, 0.22/0.7/0.10, 0.40/1.4/0.05,
    0.55/0.55/0.09, 0.68/0.9/0.07, 0.82/0.6/0.10, 0.93/0.75/0.08
  } {
    \fill[white, opacity=\o]
      ([xshift=1.4cm, yshift=-\y\paperheight]current page.north west)
      circle (\r cm);
  }
  \draw[white, line width=0.3pt, opacity=0.15]
    ([xshift=0.5cm]current page.north west) --
    ([xshift=0.5cm]current page.south west);
  \draw[white, line width=0.3pt, opacity=0.15]
    ([xshift=2.1cm]current page.north west) --
    ([xshift=2.1cm]current page.south west);
\end{tikzpicture}

% ============================================================
%  第二层：顶部/底部装饰金线带
% ============================================================
\begin{tikzpicture}[remember picture,overlay]
  \draw[GoldMain!45,line width=0.6pt]
    ([xshift=3.0cm,yshift=-1.6cm]current page.north west) --
    ([xshift=-2.0cm,yshift=-1.6cm]current page.north east);
  \foreach \x in {3.5,7,10.5,14} {
    \fill[GoldMain!55]
      ([xshift=\x cm,yshift=-1.45cm]current page.north west)
      circle (1.5pt);
  }
  \draw[GoldMain!35,line width=0.5pt]
    ([xshift=3.0cm,yshift=2.2cm]current page.south west) --
    ([xshift=-1.5cm,yshift=2.2cm]current page.south east);
  \foreach \x in {4,8,12} {
    \fill[GoldMain!40]
      ([xshift=\x cm,yshift=2.35cm]current page.south west)
      circle (1.2pt);
  }
\end{tikzpicture}

% ============================================================
%  第三层：右上角装饰几何图形
% ============================================================
\begin{tikzpicture}[remember picture,overlay]
  \draw[GoldMain!30,line width=0.7pt]
    ([xshift=-2.5cm,yshift=-2.5cm]current page.north east)
    circle (22pt);
  \draw[BlueLight!25,line width=0.5pt]
    ([xshift=-2.5cm,yshift=-2.5cm]current page.north east)
    circle (14pt);
  \fill[GoldMain!50]
    ([xshift=-2.5cm,yshift=-2.5cm]current page.north east)
    circle (2.5pt);
  \foreach \ang in {0,90,180,270} {
    \draw[GoldMain!25,line width=0.4pt]
      ([xshift=-2.5cm,yshift=-2.5cm]current page.north east)
      -- ++(\ang:28pt);
  }
  \foreach \ang in {45,135,225,315} {
    \draw[BlueLight!18,line width=0.3pt]
      ([xshift=-2.5cm,yshift=-2.5cm]current page.north east)
      -- ++(\ang:26pt);
  }
\end{tikzpicture}

% ============================================================
%  第四层：左下角数学曲线簇（ODE 相图风格）
% ============================================================
\begin{tikzpicture}[remember picture,overlay]
  \draw[GoldMain!15,line width=0.5pt,smooth]
    ([xshift=3.5cm,yshift=3.0cm]current page.south west)
    to[out=30,in=210] ++(3.0,2.0);
  \draw[BlueLight!12,line width=0.5pt,smooth]
    ([xshift=3.3cm,yshift=3.0cm]current page.south west)
    to[out=35,in=205] ++(3.2,2.2);
  \draw[GoldMain!12,line width=0.5pt,smooth]
    ([xshift=3.7cm,yshift=2.8cm]current page.south west)
    to[out=25,in=215] ++(2.8,1.8);
  \draw[BlueMain!10,line width=0.5pt,smooth]
    ([xshift=3.5cm,yshift=4.0cm]current page.south west)
    to[out=-15,in=195] ++(2.5,-1.5);
  \draw[GoldMain!10,line width=0.4pt,smooth]
    ([xshift=3.7cm,yshift=4.2cm]current page.south west)
    to[out=-20,in=190] ++(2.3,-1.3);
  \draw[GoldMain!22,-{Stealth[length=3pt,width=3pt]}]
    ([xshift=5.0cm,yshift=4.2cm]current page.south west)
    -- ++(0.8,0.5);
  \draw[BlueLight!18,-{Stealth[length=3pt,width=3pt]}]
    ([xshift=5.2cm,yshift=3.5cm]current page.south west)
    -- ++(0.7,-0.5);
\end{tikzpicture}

% ============================================================
%  第五层：公式背景区 —— 淡金色矩形 + ODE 经典公式
% ============================================================
\begin{tikzpicture}[remember picture,overlay]
  \shade[top color=GoldPale!40, bottom color=GoldPale!22]
    ([xshift=3.3cm,yshift=-3cm]current page.north west) rectangle
    ([xshift=-1.5cm,yshift=2.5cm]current page.south east);
  \draw[GoldMain!28,line width=0.4pt]
    ([xshift=3.3cm,yshift=-3cm]current page.north west) rectangle
    ([xshift=-1.5cm,yshift=2.5cm]current page.south east);

  % 递进箭头
  \draw[GoldMain!30,line width=0.4pt,-{Stealth[length=3.5pt,width=3.5pt]}]
    ([xshift=5.5cm,yshift=-4.3cm]current page.north west) --
    ([xshift=5.5cm,yshift=-4.8cm]current page.north west);
  \draw[GoldMain!30,line width=0.4pt,-{Stealth[length=3.5pt,width=3.5pt]}]
    ([xshift=5.7cm,yshift=-5.3cm]current page.north west) --
    ([xshift=5.7cm,yshift=-5.8cm]current page.north west);

  % ODE 经典公式
  \node[anchor=west,BlueMain!65] at ([xshift=3.7cm,yshift=-3.5cm]current page.north west)
    {$\displaystyle \frac{\dd y}{\dd t}=f(t,y),\quad y(t_0)=y_0$};

  \node[anchor=west,BlueLight!55] at ([xshift=3.8cm,yshift=-4.5cm]current page.north west)
    {$\displaystyle y(t)=y_0+\displaystyle\int_{t_0}^{t} f(s,y(s))\,\dd s$};

  \node[anchor=west,GoldMain!70] at ([xshift=3.9cm,yshift=-5.5cm]current page.north west)
    {$\displaystyle \dot{\vect{x}}=A\vect{x},\quad
      \vect{x}(t)=e^{tA}\vect{x}_0$};

  \node[anchor=west,BlueMain!50] at ([xshift=3.7cm,yshift=-6.5cm]current page.north west)
    {$\displaystyle \frac{\dd \vect{x}}{\dd t}
      = \vect{f}(\vect{x}),\quad
      \vect{x}\in\Omega\subset\R^n$};

  \node[anchor=west,BlueLight!45] at ([xshift=3.8cm,yshift=-7.5cm]current page.north west)
    {$\displaystyle W(t)=\det\!\bigl(\vect{x}^{(1)}(t),\dots,\vect{x}^{(n)}(t)\bigr)$};

  \node[anchor=west,GoldMain!55] at ([xshift=4.0cm,yshift=-8.5cm]current page.north west)
    {$\displaystyle \Phi(t,t_0)=\exp\!\left(\displaystyle\int_{t_0}^{t}A(s)\,\dd s\right)$};

  \foreach \y/\c in {
    -3.5/BlueMain!30, -4.5/BlueLight!28, -5.5/GoldMain!32,
    -6.5/BlueMain!25, -7.5/BlueLight!22, -8.5/GoldMain!28
  } {
    \fill[\c]
      ([xshift=-2.0cm,yshift=\y cm]current page.north west)
      circle (1.8pt);
  }
\end{tikzpicture}

% ============================================================
%  第六层：标题区 —— 使用绝对定位，避免长中文标题意外断行
% ============================================================
\begin{tikzpicture}[remember picture,overlay]
  % 系列名与短金线
  \node[
    anchor=west,
    text=BlueLight,
    font=\large\bfseries
  ] at ([xshift=4.8cm,yshift=-9.55cm]current page.north west)
    {\@seriestitle};
  \draw[GoldMain!60,line width=1.2pt]
    ([xshift=4.8cm,yshift=-10.2cm]current page.north west) --
    ([xshift=6.8cm,yshift=-10.2cm]current page.north west);
  \fill[GoldMain!70]
    ([xshift=7.1cm,yshift=-10.2cm]current page.north west) circle (2pt);
  \fill[GoldMain!45]
    ([xshift=7.55cm,yshift=-10.2cm]current page.north west) circle (1.2pt);

  % 主标题给足整行宽度，并轻微压缩字距以保证稳定排版
  \node[
    anchor=west,
    text=BlueDark,
    text width=14.2cm,
    align=left,
    inner sep=0pt,
    font=\bfseries\fontsize{28}{34}\selectfont
  ] at ([xshift=4.8cm,yshift=-11.15cm]current page.north west)
    {\sffamily\@booktitle};

  % 英文副题作为标题与署名之间的轻量过渡
  \node[
    anchor=west,
    text=GoldDark!78,
    font=\small\addfontfeatures{LetterSpace=4}
  ] at ([xshift=4.85cm,yshift=-12.5cm]current page.north west)
    {\@booksubtitle};

  % 作者与短金线
  \node[
    anchor=west,
    text=BlueMain!75,
    font=\Large
  ] at ([xshift=4.8cm,yshift=-14cm]current page.north west)
    {\@bookauthor};
  \draw[GoldMain!38,line width=0.65pt]
    ([xshift=4.8cm,yshift=-14.75cm]current page.north west) --
    ([xshift=7.5cm,yshift=-14.75cm]current page.north west);
  \fill[GoldMain!45]
    ([xshift=4.8cm,yshift=-14.75cm]current page.north west) circle (1.4pt);
  \fill[GoldMain!45]
    ([xshift=7.5cm,yshift=-14.75cm]current page.north west) circle (1.4pt);
\end{tikzpicture}

% ============================================================
%  第七层：右下角 logo
% ============================================================
\begin{tikzpicture}[remember picture, overlay]
  \begin{scope}[shift={([xshift=-5.5cm, yshift=1.5cm]current page.south east)}]
    \draw[BlueMain, line width=0.8pt] (0,0) rectangle (0.9,0.72);
    \draw[BlueMain, line width=0.8pt] (0,0.72) -- (0.45,1.08) -- (0.9,0.72);
    \draw[GoldMain, line width=0.5pt] (0.45,0.72) -- (0.45,1.08);
    \draw[BlueMain, line width=0.4pt] (0.23,0) -- (0.23,0.72);
    \draw[BlueMain, line width=0.4pt] (0.67,0) -- (0.67,0.72);
    \fill[GoldMain] (0.45,1.08) circle (1.5pt);
    \node[anchor=west, BlueMain, font=\footnotesize] at (1.1,0.32)
      {ODE Notes · Huginn};
  \end{scope}
\end{tikzpicture}

% ============================================================
%  第八层：底部封面图片
% ============================================================
\ifdefempty{\@coverimage}{}{
	\begin{tikzpicture}[remember picture,overlay]
		\node[anchor=south west,inner sep=0pt,outer sep=0pt]
			at ([xshift=3.3cm,yshift=3.05cm]current page.south west)
		{
			\includegraphics[
				width=\dimexpr\paperwidth-4.8cm\relax
			]{\@coverimage}
		};
	\end{tikzpicture}
}

% ============================================================
%  第九层：散落装饰点
% ============================================================
\begin{tikzpicture}[remember picture,overlay]
  \fill[GoldMain!18] ([xshift=-3.0cm,yshift=-9.5cm]current page.north west) circle (1.0pt);
  \fill[BlueLight!18] ([xshift=-3.8cm,yshift=-9.8cm]current page.north west) circle (0.7pt);
  \fill[GoldMain!16] ([xshift=-2.5cm,yshift=-10.2cm]current page.north west) circle (0.8pt);
  \fill[BlueMain!12] ([xshift=-4.0cm,yshift=3.5cm]current page.south west) circle (1.2pt);
  \fill[GoldMain!13] ([xshift=-5.5cm,yshift=4.0cm]current page.south west) circle (0.9pt);
  \fill[GoldMain!20] ([xshift=14.5cm,yshift=-6.0cm]current page.north west) circle (1.0pt);
  \fill[BlueLight!16] ([xshift=14.0cm,yshift=-5.5cm]current page.north west) circle (0.6pt);
  \fill[GoldMain!18] ([xshift=15.0cm,yshift=-5.8cm]current page.north west) circle (0.8pt);
\end{tikzpicture}

\end{titlepage}
}
\makeatother
